\documentclass[11pt]{article}
\usepackage[margin=1in]{geometry}
\usepackage{booktabs}
\usepackage{hyperref}
\usepackage{graphicx}
\title{What Can the Apple Neural Engine Do for Deep-Learning Training?\\
\large A Measurement Study with Energizer}
\author{Florian Grima}
\date{\today}
\begin{document}
\maketitle
\begin{abstract}
TODO.
\end{abstract}

\section{Introduction}
% Write last, once the claim is settled (docs/claims.md).

\section{Background and Related Work}
Apple's Neural Engine (ANE) is reachable only through Core~ML~\cite{coreml};
MLX~\cite{mlx2023} and PyTorch's MPS backend~\cite{paszke2019pytorch} execute on
the GPU through Metal. Apple documents ANE-oriented transformer design
choices~\cite{apple2022anetransformers}, and community reverse-engineering
efforts describe ANE behaviour beyond the public documentation
\cite{hollance2020neuralengine}.
% TODO: add reverse-engineering / private-API ANE training work after VERIFY of
% each source (see refs.bib). State clearly what we did and did not use.

\section{Methodology}
% Harness, metrics, idle-power subtraction, warm-up and repeats, thermal control:
% see docs/REPRODUCE.md.

\section{Characterisation of ANE Placement, Latency and Power}
% \input{tables/placement_ane_fraction}
% \includegraphics{figures/ane_placement.pdf}

\section{Baselines and Achieved-vs-Peak Efficiency}
% \input{tables/baselines_latency}
% \input{tables/measured_peak}

\section{Training Paths}
% MLX-native backward (train_step), hybrid overhead breakdown (hybrid_overhead),
% convergence / time-to-accuracy (convergence).

\section{Limitations}
Single chip and OS build; Core~ML placement is a planner decision that may differ
at run time; no public ANE backward pass exists, so the hybrid path pays a full
recompilation per weight refresh; results are for small models and
inference-style forward passes; no private APIs were used.

\bibliographystyle{plain}
\bibliography{refs}
\end{document}
